\section{Discusión y reflexión final (sin nota)}

\pregunta{1}{¿Qué sesgos o limitaciones identifican al promediar vectores de palabras para representar un documento completo?}
\textbf{1. Sesgos y limitaciones del promedio de vectores.} Varias, y algunas ya nos aparecieron en el Ejercicio 4. La primera es que el promedio trata todas las palabras por igual: un término clave como \code{crediticio} pesa lo mismo que un comodín como \code{requiere}, y como los comodines aparecen en todos lados, terminan arrastrando los documentos hacia el centro del espacio (así se formó el cajón de sastre del cluster 1). La segunda es que se pierde por completo el orden y la estructura de la oración: ``el banco limita el riesgo'' y ``el riesgo limita al banco'' dan exactamente el mismo vector, y una negación (``no subió la tasa'') no cambia nada la representación. Tercero, un documento que toca dos temas queda colapsado en un punto intermedio que no representa bien a ninguno; con textos más largos que estas oraciones de 8 tokens el efecto sería mucho peor. A eso se suman las limitaciones heredadas de Word2Vec: cada palabra tiene un solo vector aunque tenga varios sentidos, y el espacio arrastra los sesgos del corpus con que se entrenó (si en los textos la minería aparece siempre asociada a conflicto ambiental, por ejemplo, esa carga queda incrustada en los vectores y de ahí pasa directo al promedio).

\pregunta{2}{Comparando K-Means sobre Embeddings versus un modelo clásico como LDA o TF-IDF, ¿cuál maneja mejor los sinónimos y términos no vistos explícitamente?}
\textbf{2. K-Means sobre embeddings vs. LDA / TF-IDF.} Con los sinónimos ganan claramente los embeddings. TF-IDF (y LDA, que trabaja sobre conteos de palabras) trata cada término como una dimensión independiente: ``vivienda'' y ``casa'' son columnas distintas y dos documentos que usan una u otra no comparten nada, aunque hablen de lo mismo. En el espacio de embeddings esos sinónimos quedan cerca por haber aparecido en contextos similares, así que K-Means puede agrupar documentos temáticamente parecidos aunque no compartan ni una palabra; en nuestro corpus se nota con \code{energética}, \code{energía} y \code{eólica}, que para TF-IDF son términos sin relación alguna. Con los términos no vistos la cosa es más pareja de lo que parece: un Word2Vec entrenado desde cero tampoco puede hacer nada con una palabra fuera de su vocabulario (simplemente la descarta, como hace nuestra función \code{obtener\_document\_embedding}), igual que TF-IDF, que no tiene columna para ella. Las salidas reales son fastText, que construye el vector de una palabra nueva a partir de sus subpalabras, o usar embeddings pre-entrenados con vocabularios enormes donde es raro que un término del dominio no aparezca. O sea: embeddings mejor en sinónimos siempre, y mejor en términos no vistos solo si se usa la variante adecuada.

\pregunta{3}{¿Cómo afectaría el rendimiento del modelo el uso de embeddings contextuales (ej. BERT/BETO) frente a Word2Vec en este tipo de industria?}
\textbf{3. Embeddings contextuales (BERT/BETO) frente a Word2Vec.} Para esta industria el salto sería grande, por dos razones. La primera es la polisemia, que en el lenguaje económico-minero chileno abunda: \code{central} puede ser el banco central o una central eléctrica, \code{matriz} puede ser la matriz energética o una casa matriz, \code{banco} y \code{tasa} tienen también usos cotidianos. Word2Vec asigna un único vector por palabra, mezcla de todos sus sentidos, y eso acerca artificialmente documentos financieros y energéticos; en nuestro corpus \code{central} solo aparece junto a \code{banco}, pero en un corpus real de la industria ese cruce sería un problema constante. BERT/BETO genera un vector distinto por cada aparición según su contexto, así que ``central'' en ``banco central subió la tasa'' y en ``central de ciclo combinado'' quedarían lejos, como corresponde. La segunda razón es el tamaño del corpus, que fue la limitación transversal de todo este taller: con 12 oraciones, Word2Vec entrenado desde cero produce un espacio saturado y vecinos genéricos. BETO viene pre-entrenado sobre miles de millones de palabras en español, por lo que llega con representaciones ya robustas y solo hay que ajustarlo (o directamente usar sus embeddings de oración para el clustering del Ejercicio 4, donde probablemente los tres documentos financieros del cluster 1 habrían caído donde corresponde). El costo es computacional: inferencia mucho más pesada que una consulta a una tabla de vectores, y para vocabulario muy específico del rubro (nombres de yacimientos, jerga regulatoria) igual puede hacer falta un fine-tuning con textos del sector. Aun así, para clasificar o agrupar documentos de esta industria el intercambio conviene casi siempre.
